%% fig_pphi.tex
%% 機構の可視化: 広い断面での配向分布 P(phi)。
%% delta が小さいと phi=90deg（電場方向・最大張り出し）に単峰、
%% delta が大きいと phi1=arcsin(1/delta) 付近（軸方向）へ二峰化して移動する。
\begin{tikzpicture}
  \begin{groupplot}[
      group style={group size=2 by 1, horizontal sep=1.4cm},
      thesisaxis,
      xmin=0, xmax=360,
      xtick={0,90,180,270,360},
      xticklabels={$0$,$\tfrac{\pi}{2}$,$\pi$,$\tfrac{3\pi}{2}$,$2\pi$},
      xlabel={$\phi$}, ylabel={$P(\phi)$},
      ymin=0, ymax=0.62,
      cycle list={
        {betaA, line width=1pt},
        {betaB, line width=1pt},
        {betaC, line width=1pt},
        {betaD, line width=1pt}},
    ]
    \nextgroupplot[title={$m=3,\ \beta pE=3$}]
      % phi=90deg（小 delta の安定配向）の目印
      \draw[gray, densely dotted] (axis cs:90,0) -- (axis cs:90,0.62);
      \addplot table[x=phi_deg, y=d0333333]{\figdata/pphi_wide_m3_b3.dat};
      \addplot table[x=phi_deg, y=d1]{\figdata/pphi_wide_m3_b3.dat};
      \addplot table[x=phi_deg, y=d3]{\figdata/pphi_wide_m3_b3.dat};
      \addplot table[x=phi_deg, y=d10]{\figdata/pphi_wide_m3_b3.dat};
      \legend{$\delta=\tfrac13$,$\delta=1$,$\delta=3$,$\delta=10$}
    \nextgroupplot[title={$m=6,\ \beta pE=3$}, ylabel={}]
      \draw[gray, densely dotted] (axis cs:90,0) -- (axis cs:90,0.62);
      \addplot table[x=phi_deg, y=d0333333]{\figdata/pphi_wide_m6_b3.dat};
      \addplot table[x=phi_deg, y=d1]{\figdata/pphi_wide_m6_b3.dat};
      \addplot table[x=phi_deg, y=d3]{\figdata/pphi_wide_m6_b3.dat};
      \addplot table[x=phi_deg, y=d10]{\figdata/pphi_wide_m6_b3.dat};
  \end{groupplot}
\end{tikzpicture}
